\section{自由度与能量均分定理}
在\xref{sec:理想气体的物态方程}中讨论理想气体的压强公式和温度公式时，我们只考虑了分子的平动，更确切的说，我们将分子当成了完全弹性的质点来处理，但是这样的处理结果与实验事实并不完全相符，我们还必须要考虑分子本身的结构，单原子分子的运动只有平动，双原子分子和多原子分子则还要考虑转动和原子间的振动，因此热运动的能量也应该包含这些运动所具有的能量。为了阐明气体分子无规则运动具有的全部能量，进而求出气体内能，我们先引入气体自由度的概念。

\subsection{自由度}
确定一个物体的空间位置所需的独立坐标数，称为该物体的\uwave{自由度}。显然自由度是与物体的机械运动方式有关的物理量，我们知道，机械运动的基本方式包含：平动、转动、振动。相应的就有：\uwave{平动自由度}、\uwave{转动自由度}、\uwave{振动自由度}，三者的值我们分别用符号$t,r,s$表示。

如果物体的总自由度记为$i$，那么就有
\begin{Equation}
    i=t+r+s
\end{Equation}
现在的问题就是，如何确定自由度
\begin{itemize}
    \item 对于平动自由度，设想一个质点在三维空间自由运动，则需要三个独立坐标$x,y,z$来描述它的位置，即$t=3$，如果被限定在二维则有$t=2$，如果被限定在一维则有$t=1$。
    \item 对于转动自由度，我们知道，刚体的一般运动可以视为平动和转动的叠加，刚体平动的部分可以由质心的平动自由度给出，但是，刚体的质心确定时，刚体仍然可能有不同的空间取向，这需要引入额外的转动自由度。较一般的说，第一步我们需要确定刚体转轴的方向，这可以视为一个矢量，而矢量的方向可以由其与坐标轴的夹角$\alpha,\beta,\gamma$描述，但由于存在$\cos^2\alpha+\cos^2\beta+\cos^2\gamma=1$的关系，其中只有$\alpha,\beta$两个独立坐标。第二步我们还需要确定刚体绕转轴的角度，记为$\varphi$。因此一般情况下，我们需要三个独立坐标来描述刚体的转动，\empx{两个坐标$\alpha,\beta$给出转轴角度，一个坐标$\varphi$给出刚体绕转轴的角度}，即转动自由度$r=3$，但是如果刚体的形状本身存在一定的对称性，那么$r$可能会减小。
    \item 对于振动自由度，我们知道，刚体被认为是不会发生形变的，但是对于实际的非刚性物体而言，试想两个弹簧连接的球，其各部分间会发生相互振动，这将给出振动自由度。
\end{itemize}

\subsection{理想气体分子的自由度}
理想气体分子也可以用\xref{subsec:自由度}中定义的自由度来衡量
\begin{itemize}
    \item 单原子分子可以看作是一个能够在空间中自由运动的质点，具有$3$个平动自由度。
    \item 双原子分子具有$3$个平动自由度和$2$个转动自由度，这是因为双原子分子由两个原子通过化学键的连线构成，$3$个平动自由度确定质心，$2$个转动自由度确定化学键的角度，然而由于原子本身被视作了质点，以化学键为轴的转动是无意义的，故少了$\varphi$的自由度。
    \item 多原子分子具有$3$个平动自由度和$3$个转动自由度，这是因为多原子分子通过从不具有对称性，为自由刚体，故有全部的转动自由度，例如水分子\xce{H2O}和硫化氢分子\xce{H2S}均属折线构型，它们就具有$i=6$的自由度。然而该结论并不绝对，因为多原子分子中的各原子有时也可以排在同一直线，例如二氧化碳分子\xce{CO2}属直线构型，它的情况就与双原子分子相同，缺失了绕转轴的$\varphi$自由度，故只有$i=5$的自由度。因此自由度的判断并不能简单的依靠构成分子的原子数目来确定，还是要依据分子结构的对称性来研判。
\end{itemize}
\begin{Table}[理想气体分子的自由度]{lrrr}
<分子类型&平动自由度&转动自由度&总自由度\\>
单原子分子&$t=3$&$r=0$&$i=3$\\
双原子分子&$t=3$&$r=2$&$i=5$\\
多原子分子（原子不在同一直线上）&$t=3$&$r=3$&$i=6$\\
\end{Table}

上述考虑的都是所谓\uwave{刚性分子}的情形，即我们假定分子内的原子之间的距离保持不变，而实际分子内的原子之间还会发生振动，我们称之为\uwave{非刚性分子}，这就会导致\xref{subsec:自由度}中所述的振动自由度的出现。不过其工作方式比较特殊，在温度较低时振动自由度会被冻结，在温度升高时振动自由度才会逐渐恢复（我们可以想象一个浸水后被冰冻的弹簧，其弹性会随着温度升高冰的融化逐渐恢复）。通常来说振动自由度的“解冻温度”比较高，\empx{因此研究常温气体性质时，通常可以不考虑分子的振动自由度}。但是也有一些不能忽略振动自由度的情形，例如二氧化碳分子具有$4$个振动自由度，有论文指出在$298\si{K}$时，其振动自由度的“解冻部分”在分子平均动能上的贡献大致相当于$0.95$个自由度，这相较于$3$个平动自由度和$2$个转动自由度而言就不能被忽略了\cite{p1}。不过大学物理不考虑这种复杂情况，我们总假定分子是刚性的。

\subsection{理想气体能量按自由度均分定理}
根据\fancyref{fml:温度第二公式}和\fancyref{def:分子平均平动动能}，我们知道
\begin{Equation}*
    \xbar{\varepsilon_{\text{kt}}}=\frac{3}{2}k_BT=
    \frac{1}{2}m_0\xbar{v^2}
\end{Equation}
而由\fancyref{ppt:理想气体的速度均值性质}我们知道$\xbar{v_x^2}=\xbar{v_y^2}=\xbar{v_z^2}=(1/3)\xbar{v^2}$，这样一来
\begin{Equation}[能量按自由度均分定理初步]
    \frac{1}{2}m_0\xbar{v_x^2}=
    \frac{1}{2}m_0\xbar{v_y^2}=
    \frac{1}{2}m_0\xbar{v_z^2}=
    \frac{1}{2}k_BT
\end{Equation}
\xref{eq:能量按自由度均分定理初步}表明，分子的平均平动动能均等的分配到了每一个平动自由度上，每个自由度具有相同大小的$(1/2)k_BT$的动能。而在\xref{subsec:理想气体分子的自由度}中我们看到，分子除了平动还具有转动和振动的运动形式，为了考虑分子的全部动能，我们引入\uwave{分子平均动能} $\xbar{\varepsilon_{\text{k}}}=\xbar{\varepsilon_{\text{kt}}}+\xbar{\varepsilon_{\text{kr}}}+\xbar{\varepsilon_{\text{ks}}}$为三种运动形式产生的平均动能之和。现在的问题是，\empx{分子平均动能对于不同形式的运动，是否仍然按自由度均等分配能量}？这是个很难回答的问题，因为动能在三个平动自由度之间均分还可以从空间的各向同性去理解，但是动能在平动自由度和转动自由度等之间均等分配就没有这么显然了，这个问题的研究需要比较复杂的数学，在此略去。但事实是，这个问题有肯定的回答。
\begin{BoxTheorem}[能量按自由度均分定理]
    理想气体的分子在每一个自由度上具有相等的平均动能，其大小均为$(1/2)k_BT$。

    因此，温度为$T$时，自由度$i$的理想气体分子的平均动能为
    \begin{Equation}&[]
        \xbar{\varepsilon_{\text{k}}}=\frac{i}{2}k_BT
    \end{Equation} 
    上述结论称为\uwave{能量按自由度均分定理}。
\end{BoxTheorem}

\fancyref{thm:能量按自由度均分定理}的物理机制是，气体由非平衡态演化为平衡态的过程是依靠大量分子间无规则的频繁碰撞中交换能量试想的，在碰撞中，在一个分子上的能量可以传递给另外一个分子，在一个自由度上的能量也可以转移到另一个自由度上，在平衡态时，能量就按自由度均匀分配了，因此，所谓平衡态也可以视作分子的能量按自由度均分的标志。

\subsection{理想气体的内能}
我们知道，热现象是组成热力学系统的大量微观粒子热运动的宏观表现，因此，热现象的能量就来自微观粒子的运动和相互作用。在热学中，我们将气体分子的动能和势能的总和，称为气体的\uwave{内能}，对于理想气体，由于不考虑气体分子间的相互作用，因此气体分子间相互作用产生的势能便可以忽略不计，所以\empx{理想气体的内能就是其中所有分子无规则热运动的动能总和}。

基于上述讨论，我们就不难求出理想气体的内能了，先前的\fancyref{thm:能量按自由度均分定理}已经告诉我们单个分子上的平均动能为$(i/2)k_BT$，设想理想气体中总共有$\nu=m/M$摩尔的气体分子，而每摩尔气体分子即阿伏伽德罗常数$N_{\text{A}}$个气体分子，如果即内能为$E$
\begin{Equation}
    E=(\nu N_{\text{A}})\xbar{\varepsilon_{\text{k}}}=\frac{i}{2}\nu N_{\text{A}}k_{B}T=\frac{i}{2}\nu RT
\end{Equation}
其中最后一步\fancyref{def:玻尔兹曼常量}中应用了$R=k_{B}N_{\text{A}}$的关系。
\begin{BoxFormula}[理想气体的内能]
    理想气体的内能为
    \begin{Equation}&[]
        E=\frac{i}{2}\nu RT=\frac{i}{2}\frac{m}{M}RT
    \end{Equation}
    其中$T$为理想气体的温度，而$\nu,m,M$依次为气体的摩尔数、质量、摩尔质量。
\end{BoxFormula}
\fancyref{fml:理想气体的内能}表明，\empx{对于一定量确定种类的理想气体，其内能只和温度有关}。
这就是说理想气体的内能$E$只是温度的单值函数$R(T)$，这也与宏观的实验观测结果相符。

上述推导过程也指出，摩尔气体常量$R$和玻尔兹曼常数$k_B$分别表征了宏观和微观上的能量与温度的关系，如理想气体的分子平均动能$\xbar{\varepsilon_\text{k}}=(i/2)k_BT$，而理想气体的内能$E=(i/2)\nu RT$。

上述讨论仍有不严谨的地方，例如我们认为理想气体不存在分子间相互作用，从而就认为理想气体不存在势能了。这并不正确，因为此处的势能，\empx{不仅包含分子间势能，还包含分子内部的原子间势能}。后者可以由非刚性的理想气体分子的振动自由度给出，因为我们知道振动系统是同时具有动能和势能的，事实情况是，每个振动自由度在给出$(1/2)kT$的平均动能的同时也将产生$(1/2)kT$的平均势能，即振动自由度对内能的贡献是同等量的平动自由度和转动自由度的两倍，以二氧化碳为例的话，当温度足够高使其四个振动自由度完全解冻时，其内能将是$E=(13/2)\nu RT$（$13=3+2+4\times 2$），而非\xrefpeq{}的$E=(9/2)\nu kT$（$9=3+2+4$）。

